\documentclass{article}
\usepackage[utf8]{inputenc}
\usepackage{amsmath, amssymb, amsthm, booktabs}

\title{Comprehensive Reference: Hyperkähler Manifolds}
\author{}
\date{}

\begin{document}

\maketitle

\section{1. Algebraic Foundations}
A \textbf{Hyperkähler manifold} $(M, g)$ of real dimension $4n$ is equipped with a triplet of complex structures $(I, J, K)$ acting on the tangent bundle. These satisfy the quaternionic algebra:
\[ I^2 = J^2 = K^2 = IJK = -1 \]
These structures are compatible with the Riemannian metric $g$, satisfying:
\[ g(IX, IY) = g(JX, JY) = g(KX, KY) = g(X, Y) \]

\section{2. Parallelism and Holonomy}
The manifold is Hyperkähler if and only if $\nabla I = \nabla J = \nabla K = 0$. This implies the three Kähler forms are parallel, restricting the holonomy group $\text{Hol}(g)$ to a subgroup of $Sp(n)$, forcing Ricci-flatness:
\[ \text{Ric}(g) = 0 \]

\section{3. Twistor Space Construction}
For a Hyperkähler manifold $M$, we define the \textbf{Twistor space} $Z = M \times \mathbb{CP}^1$. The complex structures $I, J, K$ define a bundle of complex structures over $\mathbb{CP}^1$. A point $(\zeta_1 : \zeta_2) \in \mathbb{CP}^1$ corresponds to a complex structure $\mathcal{I}$ on $M$:
\[ \mathcal{I} = \frac{\zeta_1 I + \zeta_2 J + \zeta_3 K}{|\zeta|} \]


\section{4. The Hyperkähler Quotient}
Let $G$ act on $M$ preserving the triplet $(I, J, K)$ with hyper-moment map $\mu = (\mu_I, \mu_J, \mu_K)$. The quotient is:
\[ M /// G = \mu^{-1}(0, 0, 0) / G \]

\section{5. Gibbons-Hawking and ALE Classification}
The Gibbons-Hawking construction defines the metric via a harmonic function $V$ on $\mathbb{R}^3$:
\[ ds^2 = V^{-1}(d\tau + \vec{\omega} \cdot d\vec{x})^2 + V d\vec{x}^2 \]

ALE (Asymptotically Locally Euclidean) spaces are classified by the minimal resolution of Kleinian singularities $\mathbb{C}^2/\Gamma$, corresponding to the $ADE$ Dynkin diagrams.


\section{6. Manifold Hierarchy}
\begin{table}[h!]
\centering
\begin{tabular}{lccc}
\toprule
\textbf{Manifold} & \textbf{Complex Struct.} & \textbf{Holonomy} & \textbf{Ricci-Flat} \\
\midrule
Kähler & 1 & $U(n)$ & No \\
Calabi-Yau & 1 & $SU(n)$ & Yes \\
Hyperkähler & 3 & $Sp(n)$ & Yes \\
\bottomrule
\end{tabular}
\end{table}

\section*{References}
\begin{thebibliography}{99}
\bibitem{Calabi79} Calabi, E. (1979). Métriques kählériennes et fibrés holomorphes. \emph{Ann. Sci. École Norm. Sup.}
\bibitem{Beauville83} Beauville, A. (1983). Variétés Kähleriennes dont la première classe de Chern est nulle. \emph{J. Differential Geom.}
\bibitem{Berger55} Berger, M. (1955). Sur les groupes d'holonomie. \emph{Bull. Soc. Math. France}
\bibitem{HKLR87} Hitchin, N. J., et al. (1987). Hyperkähler metrics and supersymmetry. \emph{Comm. Math. Phys.}
\bibitem{AHS78} Atiyah, M. F., et al. (1978). Self-duality in four-dimensional Riemannian geometry. \emph{Proc. Roy. Soc. Lond.}
\bibitem{GH78} Gibbons, G. W., \& Hawking, S. W. (1978). Classification of gravitational instanton symmetries. \emph{Phys. Lett. B}
\bibitem{Kronheimer89} Kronheimer, P. B. (1989). The construction of ALE spaces as hyperkähler quotients. \emph{J. Differential Geom.}
\end{thebibliography}

\end{document}
